\section{Contractions and the cyclic insertion map}

Throughout, $\lVert\cdot\rVert$ is the operator norm of a complex matrix acting
between Euclidean spaces, and a contraction is a matrix or operator of norm at
most one.

\begin{lemma}[Elementary contractions]
    \label{lem:matrix_l2_contractions}
    \lean{Matrix.l2_opNorm_le_one_of_conjTranspose_mul_self_le_one,
        Matrix.IsIsometry.l2_opNorm_le_one,
        Matrix.l2_opNorm_conjTranspose_mul_mul_le_one,
        Matrix.l2_opNorm_list_prod_le_one,
        Matrix.l2_opNorm_kronecker_le_one,
        Matrix.l2_opNorm_toMatrix_toPEquiv_le,
        Matrix.norm_submatrix_equiv_le}
    \leanok
    Let $A$ and $F$ be complex matrices indexed by finite sets.
    \begin{enumerate}
        \item If $\lVert A^\dagger A\rVert\le1$, then $\lVert A\rVert\le1$; in
              particular every isometry $A^\dagger A=I$ is a contraction.
        \item If $A$ and the square matrix $F$ are contractions, so is
              $A^\dagger FA$.
        \item A finite product of square contractions is a contraction, and
              the Kronecker product of two contractions is a contraction.
        \item The matrix of a bijection between two finite coordinate sets is a
              contraction, and relabelling both coordinates of a square matrix
              along one bijection does not increase its norm.
    \end{enumerate}
\end{lemma}

\begin{proof}\leanok
    The identity $\lVert A^\dagger A\rVert=\lVert A\rVert^2$ gives the first
    statement. The second and third follow from submultiplicativity, the
    identity $\lVert A^\dagger\rVert=\lVert A\rVert$, induction on the length
    of the product, and $\lVert A\otimes B\rVert\le\lVert A\rVert\,\lVert B\rVert$.
    The matrix of a bijection is a relabelling of the identity matrix, and
    relabelling coordinates is conjugation by isometries.
\end{proof}

\begin{lemma}[Operator norm of a block-diagonal matrix]
    \label{lem:block_diagonal_operator_norm}
    \lean{Matrix.l2_opNorm_le_of_blocks}
    \leanok
    Let $Y$ be a complex matrix with rows indexed by a finite set $\alpha$ and
    columns indexed by a finite set $\beta$, and let $u:\alpha\to U$ and
    $u':\beta\to U$ be labellings by a finite set $U$ such that $Y_{ab}=0$
    whenever $u(a)\ne u'(b)$. If $c\ge0$ and, for every label $z\in U$,
    \begin{align}
        \sum_{u(a)=z}\ \sum_{u'(b)=z}\lvert Y_{ab}\rvert^2\le c^2,
        \notag
    \end{align}
    then $\lVert Y\rVert\le c$.
\end{lemma}

\begin{proof}\leanok
    For a vector $v$, the entry $(Yv)_a$ involves only the coordinates $v_b$
    with $u'(b)=u(a)$, so the Cauchy--Schwarz inequality bounds
    $\lvert(Yv)_a\rvert^2$ by the squared Hilbert--Schmidt norm of the row
    restricted to that block times the squared norm of the corresponding part
    of $v$. Summing over $a$, grouped by labels, gives
    $\lVert Yv\rVert^2\le c^2\lVert v\rVert^2$.
\end{proof}

\begin{lemma}[Regrouping Kronecker products]
    \label{lem:kronecker_shuffle}
    \lean{Matrix.kronecker_kronecker_eq_submatrix,
        Matrix.toMatrix_toPEquiv_prodCongr,
        Matrix.toMatrix_toPEquiv_submatrix}
    \leanok
    For complex matrices $A_1,B_1,A_2,B_2$,
    \begin{align}
        (A_1\otimes B_1)\otimes(A_2\otimes B_2)
        =\Pi\,\bigl((A_1\otimes A_2)\otimes(B_1\otimes B_2)\bigr)\,\Pi'^{\dagger},
        \notag
    \end{align}
    where $\Pi$ and $\Pi'$ exchange the two middle factors of the row and
    column coordinates. The matrix of a product $f\times g$ of bijections is the
    Kronecker product of the matrices of $f$ and $g$, and relabelling the
    coordinates of the matrix of a bijection gives the matrix of the composed
    bijection.
\end{lemma}

\begin{proof}\leanok
    Compare entries: both sides of the first identity have entry
    $(A_1)_{aa'}(B_1)_{bb'}(A_2)_{cc'}(B_2)_{dd'}$ at the corresponding
    coordinates, and the matrix of a bijection has entry one exactly on its
    graph.
\end{proof}

\begin{lemma}[Product vectors]
    \label{lem:kronecker_vectors}
    \lean{EuclideanSpace.vecKron,
        EuclideanSpace.norm_vecKron,
        EuclideanSpace.norm_vecKron_self_sub_le,
        EuclideanSpace.norm_toLp_comp_equiv}
    \leanok
    For vectors $x,x'\in\C^m$ and $y\in\C^n$, the product vector
    $(x\otimes y)_{(i,j)}=x_iy_j$ satisfies
    $\lVert x\otimes y\rVert=\lVert x\rVert\,\lVert y\rVert$ and
    \begin{align}
        \lVert x\otimes x-x'\otimes x'\rVert
        \le\lVert x-x'\rVert\bigl(\lVert x\rVert+\lVert x'\rVert\bigr).
        \notag
    \end{align}
    Relabelling the coordinates of a vector along a bijection preserves its
    norm.
\end{lemma}

\begin{proof}\leanok
    The squared norm of $x\otimes y$ is the product of the squared norms.
    Writing $x\otimes x-x'\otimes x'=(x-x')\otimes x+x'\otimes(x-x')$ and using
    the triangle inequality gives the estimate.
\end{proof}

\begin{lemma}[Matrices acting on Euclidean space]
    \label{lem:matrix_euclidean_action}
    \lean{Matrix.inner_toEuclideanCLM_eq_dotProduct,
        Matrix.norm_toEuclideanCLM_le_of_norm_le_one,
        Matrix.toEuclideanLin_one_apply,
        Matrix.toEuclideanLin_mul_apply,
        Matrix.schmidtCoeffMatrix_kronecker_one_mulVec}
    \leanok
    For a square complex matrix $A$ and vectors $x,y$, the inner product
    $\langle x,Ay\rangle$ is the sum $\sum_i\overline{x_i}(Ay)_i$, and
    $\lVert Ax\rVert\le\lVert x\rVert$ when $\lVert A\rVert\le1$. The identity
    matrix acts as the identity and a product $AB$ acts as the composition.
    For a vector $\psi$ on a product $\alpha\times\beta$ with coefficient
    matrix $\Psi_{ab}=\psi_{(a,b)}$ and a matrix $K$ with columns indexed by
    $\alpha$, the vector $(K\otimes I)\psi$ has coefficient matrix $K\Psi$.
\end{lemma}

\begin{proof}\leanok
    Each statement is a direct computation in coordinates.
\end{proof}

\begin{lemma}[Rescaling an approximate contraction]
    \label{lem:rescaled_contraction}
    \lean{norm_le_one_add_of_norm_sub_le,
        norm_inv_one_add_smul_le_one,
        norm_inv_one_add_smul_sub_le}
    \leanok
    Let $\delta\ge0$ and let $A,B$ be elements of a complex normed space with
    $\lVert B\rVert\le1$ and $\lVert A-B\rVert\le\delta$. Then
    $\lVert A\rVert\le1+\delta$, the rescaling $(1+\delta)^{-1}A$ has norm at
    most one, and
    \begin{align}
        \lVert(1+\delta)^{-1}A-B\rVert\le2\delta.
        \notag
    \end{align}
\end{lemma}

\begin{proof}\leanok
    The first bound is the triangle inequality. Since
    $\lVert(1+\delta)^{-1}A-A\rVert=\delta(1+\delta)^{-1}\lVert A\rVert\le\delta$,
    the triangle inequality gives the last estimate.
\end{proof}

Let $\eta$ be a unit vector of a finite-dimensional Hilbert space $\mathcal W$
and let $m\ge1$.

\begin{definition}[Cyclic insertion map]
    \label{def:peps_cyclic_insertion}
    \lean{CyclicInsertion.insertAt, CyclicInsertion.cyclicInsertion,
        CyclicInsertion.idealInsertion}
    \leanok
    For $k\in\{0,\dots,m-1\}$, let $I_k:\mathcal W\to\mathcal W^{\otimes m}$
    place its input in factor $k$ and $\eta$ in the remaining $m-1$ factors.
    Define the cyclic and ideal insertion maps by
    \begin{align}
        T_m & =\frac1m\sum_{k=0}^{m-1} I_k,          &
        P_m & =|\eta^{\otimes m}\rangle\langle\eta|.
        \label{eq:peps_cyclic_insertion}
    \end{align}
\end{definition}

\begin{lemma}[Insertions as permutations of the factors]
    \label{thm:peps_insertion_permutation}
    \lean{CyclicInsertion.insertAt_perm,
        CyclicInsertion.cyclicInsertion_eq_sum_addRight,
        CyclicInsertion.registerPerm_pair_index}
    \leanok
    \uses{def:peps_cyclic_insertion}
    For any permutation $\sigma$ of $\{0,\dots,m-1\}$ with induced factor
    permutation $U_\sigma$ on $\mathcal W^{\otimes m}$, $I_{\sigma(k)}=U_\sigma I_k$,
    so $T_m=m^{-1}\sum_{k=0}^{m-1}U_{(+k)}I_0$. On
    $\mathcal W=\mathcal U\otimes\mathcal V$, $U_\sigma$ permutes the
    $\mathcal U$ and $\mathcal V$ coordinates by $\sigma$.
\end{lemma}

\begin{proof}\leanok
    Verify $I_{\sigma(k)}=U_\sigma I_k$ and the pair-coordinate action of
    $U_\sigma$ on product basis vectors.
\end{proof}

\begin{lemma}[Contractivity and the fixed vector]
    \label{thm:peps_cyclic_insertion_basic}
    \lean{CyclicInsertion.norm_cyclicInsertion_le_one,
        CyclicInsertion.cyclicInsertion_self}
    \leanok
    \uses{def:peps_cyclic_insertion}
    The map $T_m$ satisfies $\|T_m\|\le 1$ and $T_m\eta=\eta^{\otimes m}$.
\end{lemma}

\begin{proof}\leanok
    Since $\|\eta\|=1$, each $I_k$ is an isometry with $I_k\eta=\eta^{\otimes m}$,
    so~(\ref{eq:peps_cyclic_insertion}) gives
    $\|T_m\|\le m^{-1}\sum_{k=0}^{m-1}\|I_k\|=1$ and $T_m\eta=\eta^{\otimes m}$.
\end{proof}

\begin{lemma}[One-effect estimate]
    \label{thm:peps_one_effect_estimate}
    \lean{CyclicInsertion.norm_cyclicInsertion_of_inner_eq_zero,
        CyclicInsertion.norm_cyclicInsertion_sub_ideal_le,
        CyclicInsertion.norm_cyclicInsertion_sub_ideal_eq,
        CyclicInsertion.cyclicInsertion_sub_ideal_eq_zero,
        CyclicInsertion.norm_rTensor_cyclicInsertion_sub_ideal_le}
    \leanok
    \uses{def:peps_cyclic_insertion}
    For every $v\perp\eta$, $\|T_mv\|=\|v\|/\sqrt m$. Hence
    $\|T_m-P_m\|\le m^{-1/2}$ (with equality if $\{\eta\}^\perp\neq\{0\}$ and
    $T_m=P_m$ if $\{\eta\}^\perp=\{0\}$), and
    $\|(T_m-P_m)\otimes\mathbf 1_{\mathcal S}\|\le m^{-1/2}$ for any spectator
    space $\mathcal S$.
\end{lemma}

\begin{proof}\leanok
    \uses{thm:peps_cyclic_insertion_basic}
    For $v\perp\eta$ and $j\neq k$,
    $\langle I_jv,I_kv\rangle=|\langle v,\eta\rangle|^2\|\eta\|^{2(m-2)}=0$
    and $\|I_kv\|=\|v\|$, so $\|\sum_{k=0}^{m-1}I_kv\|^2=m\|v\|^2$ and
    $\|T_mv\|=\|v\|/\sqrt m$. For any $v\in\mathcal W$, set
    $v_\perp=v-\langle\eta,v\rangle\eta\perp\eta$; by
    Lemma~\ref{thm:peps_cyclic_insertion_basic},
    $\|(T_m-P_m)v\|=\|T_mv_\perp\|=\|v_\perp\|/\sqrt m\le\|v\|/\sqrt m$, with
    equality on any nonzero $v\in\{\eta\}^\perp$. Tensoring with
    $\mathbf 1_{\mathcal S}$ preserves the norm bound.
\end{proof}

\begin{lemma}[Combining pair vectors]
    \label{thm:peps_pair_combine}
    \lean{EuclideanSpace.pairCombine, EuclideanSpace.norm_pairCombine}
    \leanok
    For $\eta\in\mathcal U\otimes\mathcal V$ and
    $\eta'\in\mathcal U'\otimes\mathcal V'$, the regrouped pair vector
    $\eta\otimes\eta'\in(\mathcal U\otimes\mathcal U')\otimes(\mathcal V\otimes\mathcal V')$
    satisfies $\|\eta\otimes\eta'\|=\|\eta\|\,\|\eta'\|$.
\end{lemma}

\begin{proof}\leanok
    Regrouping the product coordinates factors
    $\|\eta\otimes\eta'\|^2=\|\eta\|^2\|\eta'\|^2$.
\end{proof}
